
\section{Formale Aspekte , dieses Kapitel kommt in der Arbeit natürlich raus}
%
\subsection{Mathematische Gleichungen}
Eine mehrzeilige Gleichung sieht so aus (die Symbole nach den und-Zeichen werden untereinander gesetzt). Die nonmber-Befehle verhindern dass die Gleichung nummertiert wird (Geschmackssache, ist nie falsch wenn eine Gleichung nummeriert ist). Aber: eine Gleichung auf die man refernziert (also die ein Label hat), muss nummeriert sein!
\begin{align}
    A &= \sum_{i=1}^N x_i \label{eq:1}\nonumber\\
    B &= \frac{\pi}{2}
\end{align}

Eine inline-Gleichung: $x=45b + \frac{2}{3}\pi$. Der Text geht weiter! Auf inline-Gleichungen kann man keine Refernzen erstellen.

\subsection{Das ist eine Auflistung}
\begin{enumerate}
\item Element 1
\item Element 2
\end{enumerate}

\subsection{Das ist eine Bullet-Liste}
\begin{itemize}
\item Element 1
\item Element 2
\end{itemize}

\subsection{Eine Grafik bindet man so ein}
Zulässige Formate sind generell eps, pdf und png. Sie referenzieren eine Abbildung als Abb. \ref{fig:bildchen}.
\begin{figure}[t!]
    \centering
    \includegraphics[width=0.8\textwidth]{logo.pdf}
    \caption{Logo der HAW Fulda}
    \label{fig:bildchen}
\end{figure}
\begin{figure}[b!]
	\centering
	\begin{subfigure}[b]{6cm}
	  \centering
    %\fbox{\parbox{5cm}{\centering ~\vspace{1.5cm}\\Dummy\\~\vspace{1.5cm}}} 
    \includegraphics[width=5cm]{logo.pdf}
		\caption{Caption of subfigure a (can be empty)}
		\label{fig:1}
	\end{subfigure}
	\begin{subfigure}[b]{6cm}
    \centering
    %\fbox{\parbox{5cm}{\centering ~\vspace{1.5cm}\\Dummy\\~\vspace{1.5cm}}} 
    \includegraphics[width=5cm]{logo.pdf}
		\caption{Caption of subfigure b (can be empty)}
		\label{fig:2}
	\end{subfigure}
\caption{Figure using subfigures}
\label{fig:figure_with_subfigures}
\end{figure}
Eine komplizierte Abbildung mit Unterabbildungen ist auch möglich, wird referenziert als Abb. \ref{fig:figure_with_subfigures}.


\subsection{Code-Listings}
Es gibt zwei Möglichkeiten, code einzubinden: entweder als Grafik, typischerweise im PNG-Format, oder als \verb|lstlisting| environment. Ein Beispiel sehen Sie in Listing~\ref{lst:listing}.
\begin{lstlisting}[label=lst:listing, caption={Ein Python-Listing. Syntax-Highlighting uvm. ist möglich, prüfen Sie die Optionen des listing-Pakets.}]
  # Factorial in PYthon
  def f(x):
    if x == 1:
      return 1
    else:
      return x*f(x-1)
\end{lstlisting}

\subsection{So schreibt man einen Algorithmus}

\begin{algorithm}[H]
 \KwData{this text}
 \KwResult{how to write algorithm }
 initialization\;
 \While{not at end of this document}{
  read current\;
  \eIf{understand}{
   go to next section\;
   current section becomes this one\;
   }{
   go back to the beginning of current section\;
  }
 }
 \caption{How to write algorithms\label{alg:dummy}
 }
\end{algorithm}

\subsection{So gestaltet man eine Tabelle}

\begin{table}[H]
\caption{Beispielstabelle\label{tab:beispiel}
}
\centering
\begin{tabular}{llr}
\hline
A    & B & C \\
\hline
D      & per gram    & 11.65      \\
          & each        & 1.01       \\
E       & stuffed     & 32.54      \\
F       & stuffed     & 73.23      \\
G & frozen      & 8.39       \\
\hline
\end{tabular}
\end{table}


\subsection{Interne Referenzen}
So wird ein Kapitel oder Unterkapitel referenziert: Kap.~\ref{sec:einleitung},
Kap.~\ref{sec:webquellen}. Auf Gleichungen bezieht man sich so: Wie in Gl.~(\ref{eq:1}) gezeigt,
sehen Gleichungen in der Regel gut aus. Auf Abb.~\ref{fig:bildchen} bezieht man sich so. Auf
Tab.~\ref{tab:beispiel} referenziert man so. Algorithmen sind analog: siehe Alg.~\ref{alg:dummy}.
Generell kann man alles zitieren was ein Label hat.

\subsection{Textformatierung}
\textbf{So wird dick geschrieben} und \textit{so kursiv}.

\subsection{Zitieren}\label{sec:zitate}
Generell werden indirekte Zitate so benutzt: wie in \cite{clemen1989combining} gezeigt, gilt blablaba...
Direkte Zitate können so aussehen: In \cite{gilabert2006intelligent} wird folgendes festgestellt: "Blablabla". 
Für jedes zitierte Werk ist ein BibTex-Eintrag nötig! Eine gute Quelle ist Google Scholar. Nutzen Sie darüber hinaus die Angebote der Bibliotheken sowie geeignete Online-Quellen wie z.B.:
\begin{itemize}
\item ACM Digital Library: http://www.acm.org/dl
\item IEEE: http://www.computer.org
\item  CiteSeer: http://citeseer.ist.psu.edu
\end{itemize}
Mehr zu direkten und indirekten Zitaten finden Sie in \ref{app:zitieren}.

\subsection{Webquellen zitieren}\label{sec:webquellen}
So wird eine URL/Webquelle zitiert: \cite{shiny1}, siehe auch den Eintag im BibTeX-File.
Wichtig: für jede Web-Quelle ein BibTeX-Eintrag! Wenn Sie das auf die hier gezeigte Art machen, werden URLs (fast) automatisch getrennt. Kontrollieren Sie trotzdem die Literaturliste, es kann sein dass das nicht immer funktioniert. Bie Webquellen muss unbedingt das Datum das Abrufs angegeben werden! Am besten benutzt man dafür den "note"-Eintrag in BibTeX.

\subsection{Literaturverzeichnis erstellen}
Hierzu müssen BibTeX-Einträge in die Datei literatur.bib eingefügt werden. Die BibTeX-Keys sind jeweils Argumente für die cite-Kommandos! Falls Sie lokal (also nicht in Overleaf) arbeiten: Wenn Sie literatur.bib ändern müssen Sie alles mindestens 5x compilieren: 3x mit latex, 1x mit BibTex und dann noch 2x mit LaTeX (in der Reihengfolge). Am besten Sie machen ein Skript dafür! Overleaf macht das natürlich automatisch...

Zitierbare Quellen (wissenschaftliche Zeitschriften oder Journale mit DOI) sind Internetquellen (Blogs, Webseiten, ...) vorzuziehen. Prüfen Sie besonders im Netz kritisch die Qualität der Quelle! Wenn Sie eine Internetquelle zitieren, geben Sie stets das Datum des Abrufs an, damit man sie mit Hilfe von Archiv-Seiten auffinden kann. Typische bibliografische Angaben bei Quellen sind: 
Name, Jahr, Autor*innen, Konferenz/Journal, ggf. Seitenangabe, Nummer/Volume, Abrufdatum bei Webseiten. Seitenangaben sind nur nötig wenn man einen spezifischen Teil einer Arbeit zitiert.

%\subsection{Erstellung eines PDFs im PDF-A Format}
%Durch Einbinden geeigneter Packages (pdfx) wird diese Vorlage bereits als PDF-A erzeugt. Sie sollten allerdings die Metadaten in der Datei main.xmpdata %anpassen!

\section{Code Snippets}
\begin{figure}
    \centering
    \includegraphics[width=0.5\textwidth]{snippet.png}
    \caption{Normalerweise bindet man Snippets als Bilder ein. Falls sie mehr als eine Viertelseite umfassen, bitte in den Anhang so wie hier gezeigt. Diese Snippets können dann trotzdem aus dem Text referenziert werden.\label{Snippet}
    }
\end{figure}
Binden Sie Snippets ein wie in \ref{Snippet} gezeigt!

\section {Generelle Hinweise}
\begin{itemize}
    \item Diese Vorlage ist nicht universell! Jede wissenschaftliche Arbeit ist anders. Daher: Halten Sie Kontakt zu Ihren Betreuer*innen und reden Sie mit ihnen über die konkrete Ausformung der Arbeit
    \item Sie müssen nicht LaTeX benutzen, Word ist auch ok. In jedem Fall sollte die Form dieser Vorlage nachempfunden sein (Schriftgröße, Zeilenabstand, ...).
    \item Bei Fragen sind stets die Betreuer*innen die erste Anlaufstation, bitte nutzen Sie diese Möglichkeit, v.a. in den Sprechstunden!
    \item Generell wird das Gendern weder vorgeschrieben noch verboten, Sie können dies nach eigenem Ermessen gestalten. 
    \item Sprache: Deutsch oder Englisch
     \item Die Länge einer Arbeit kann variieren, je nach Thema. Weniger als 40 Seiten wäre allerdings ungewöhnlich!
     \item Besonderes Augenmerk sollte auf jeden Fall auf einer ansprechenden Form liegen (schöne und ausreichend viele Grafiken wo nötig, Formatierung, etc.) sowie der Sprache (keine Umgangssprache, elegantes Formulieren, Benutzung von Nebensätzen, ...) liegen, da diese beiden Punkte stark in die Note eingehen. Ebenfalls sehr wichtig (für die Benotung, aber auch so) sind die Kapitel Einleitung, Diskussion und Schluss, da oft nur diese gelesen werden. KI-Tools wie Grammarly können benutzt werden um die Sprache zu verbessern, ebenso wie Rechtschreib- und Grammatik-Checker (z.B. aspell). 
     \item Achten Sie bei Grafiken darauf, dass sie a) lesbar sind wenn das Dokument in A4-Größe betrachtet wird b) nicht verpixelt sind.
     \item Die Ich-Form ist zu vermeiden. Also 'es wurde festgestellt...' statt 'ich habe herausgefunden...'. Generell keine Prosa und chronologische Darstellung ('zuerst wurde das gemacht, dann das und dann das') sondern inhaltliche Gliederung!
     \item Abkürzungen, sofern sie nicht absolut offensichtlich sind, sollten bei der ersten Benutzung eingeführt werden. Ein Abkürzungsverzeichnis kann sinnvoll sein, vor der Einleitung.
     \item Achten Sie darauf, Abschnitte nicht zu sehr zu verschachteln. Maximal zwei Unterebenen ist eine gute Regel (also 1.1.1 aber nicht 1.1.1.2). 
\end{itemize}

\section{Mehr zu direkten und indirekten Zitaten}\label{app:zitieren}
Zentral für das wissenschaftliche Arbeiten ist, dass man den Ursprung jeder Darstellung, sei es nun
eine Tatsache oder eine Bewertung, eindeutig kenntlich macht: \cite{shiny1}. Diese
Forderung hat drei Gründe:
\begin{itemize}
\item Man darf sich nicht „mit fremden Federn schmücken“, d.h. fremde Gedanken als die
eigenen ausgeben. Dies wäre der Fall, wenn man übernommene bzw. referierte Aussage
nicht als solche kenntlich macht.
\item  Man kann so die Lesenden auf weitere Literatur zur Vertiefung verweisen, wenn man aus
Platz- und Strukturgründen nicht selbst alle Details eines Gedankens oder Wissensbau-
steins referieren möchte. Dadurch kann man sich auf das Wesentliche beschränken,
welches zum Verständnis und wissenschaftlichen Einordnen der eigenen Arbeit wichtig ist.
\item Durch den genauen Nachweis oder Beleg wird sichergestellt, dass sich keine Fehler in
Form von falschen Zitaten in der Wissenschaft festsetzen können. Jeder kann durch den
Beleg genau die referierte Stelle im Original nachlesen und feststellen, ob der Autor richtig
verstanden und wiedergegeben wurde.
\item Man schützt sich durch das Belegverfahren davor, für fremde Fehler „gerade stehen zu
müssen“. D.h. wenn einem Autor ein Fehler unterläuft, den Sie in Ihrer Bachelor-
/Masterarbeit ohne Beleg wiedergeben, gilt das so, als ob Sie selbst diesen Fehler
gemacht haben. Wenn Sie aber belegen, wo Sie eine falsche Aussage, die Sie nicht als
falsch erkannt haben, „abgeschrieben“ haben, dann liegt die Verantwortung für den Inhalt
nicht bei Ihnen, sondern bei Ihrer Quelle.
\end{itemize}
Man unterscheidet zwischen zwei Arten von wissenschaftlichen Belegen, dem Verweis und dem
Zitat.
\subsection{Verweis oder indirektes Zitat}
Ein Verweis bezeichnet die sinngemäße Wiedergabe eines Sachverhalts in eigenen Worten. Nach
den entsprechenden Ausführungen wird der Verweis \cite{clemen1989combining} eingefügt. 
Eine gute Möglichkeit kenntlich zu machen, dass man sich auf einen bestimmten Autor bezieht, ist
die Verwendung solcher Formulierungen wie 'nach Darstellung in \cite{clemen1989combining} ...' oder 'Wie in \cite{clemen1989combining} ausgeführt, ...'.
\subsection{Zitat}
Ein Zitat bezeichnet die exakte wörtliche Wiedergabe aus einem Text. Ein Zitat muss immer in
Anführungszeichen stehen und mit Angabe der Quelle (genauso wie beim Verweis). Dies ist in der Informatik eher unüblich und sollte sehr sparsam eingesetzt werden.

\section{Hinweise zum Kolloquium}
Das Kolloquium besteht aus 15 Minuten Präsentation gefolgt von ca 20 Minuten Fragen zur Arbeit und Diskussion. Anwesend sind mindestens die beiden Gutachter*innen, Dritte können ohne weitere Anmeldung ebenfalls dabei sein (Kolloquien sind öffentlich). 

Zielgruppe sollen Laien (z.B. höheres Management) sein, nicht die beiden Gutachter*innen. Das bedeutet dass mindestens ein Drittel der Präsentation aus Motivation und Kontext besteht. Hier kann sich stark an die Einleitung der Arbeit angelehnt werden. Viele Bilder helfen ebenfalls, können aus der Arbeit sein.

Maximal soll die Präsentation 15 Slides enthalten (also maximal 1 Slide pro Minute). Es können zusätzliche Slides mit Details erstellt werden die vermutlich in der Diskussion sowieso gefragt werden.

Das Kolloquium ist nicht benotet, dh. es gibt nur bestanden oder nicht bestanden als Ergebnis. 

Generell sind die Studierenden verantwortlich für die Terminfindung und sollen nach Einreichung der Arbeit selbst aktiv werden. 

Tipps:
\begin{itemize}
    \item langsam sprechen! Laut sprechen!
    \item Unter Zeitdruck lieber ganze Folien weglassen als schneller sprechen
    \item Folien nummerieren
    \item Bitte unbedingt 15min vor dem Termin Hardware (Beamer, Headset, Kamera) testen. Das macht einen extrem schlechten Eindruck wenn dadurch Verzögerung entsteht
    \item niemals nur den Text der Folien vorlesen. Folien sollen sehr wenig Text und viele Bilder haben. Text soll nur Stichpunkte angeben, der Rest wird frei formuliert
    \item Kolloquien können ohne Weiteres online stattfinden!
    \item bei physischen Terminen: immer einen USB-Stick mit der Präsentation als PDF dabeihaben falls man einen anderen PC nutzen muss. PDF geht überall, PowerPoint nicht (und sieht v.a. nicht überall gleich aus)
    \item Totales No-Go: Rechtschreibfehler auf Folien
    \item Vermeiden Sie aufwändige Animationen oder Videos wenn möglich, kann immer sein dass das nicht funktioniert
\end{itemize}
